\section{Consciousness}

The cleanest way to arrive at consciousness is to reach it without using the word on the way up. The ladder from
field to human, built rung by rung, never needed the term, and that is the point. Consciousness in this theory is
not a special extra that switches on once a system is complicated enough. It is integrated experience, and both of
its ingredients are already in the model. Before going further, one thing must be stated plainly. The footing of
this chapter is different from the footing of the physics. The physics is measured. Consciousness here rests on a
premise, and that premise is marked as such throughout. This part of the theory is interpretive and speculative,
and it should be read in that spirit.

\subsection{The sharp distinction}

Two notions are kept rigorously apart. \emph{Intelligence} is problem-solving capability. \emph{Consciousness} is
integrated experience, the felt going-through of reality. A system can in principle have one without the other, a
fast solver with little integration and little felt unity, or a richly integrated whole that solves nothing in
particular. Blurring the two is the most common error in this area, and the theory refuses it from the start. The
previous chapter handled intelligence. This one handles experience.

\begin{definition}[Consciousness]
\emph{Consciousness} is integrated experience. Its amount is set by how much a region of experience is bound into
a single whole. It is distinct from intelligence, which is problem-solving capability, and the two can occur
apart.
\end{definition}

\subsection{The two ingredients}

The first ingredient is experience, and it enters as a premise rather than a result. This is the founding axiom
already stated at the start of the paper, that experience exists and the one substance is the vibe, a felt state.
Experience is not produced by matter here. It is the ground, and matter is a derived face of it. The experiments
do not prove that the substrate feels. They assume it and then build, and the candor of the whole chapter depends
on saying so without hedging.

The second ingredient is integration, and this one is measured. Integration is how much a region is bound into a
single thing, with more internal coupling than external coupling, and it is the same quantity that picks out
selves and around which a self-model forms \cite{pollard2026vibetest}. With the premise and the measure in hand the
synthesis is short. Consciousness is what integration does to experience. If experience is the ground and
integration binds it, then consciousness is integrated experience, and its amount is fixed by how much integration
there is. No new ingredient is introduced.

\begin{proposition}[Integration is what makes experience unified]
Given experience as the ground and integration as the measure of how much a region is bound into one, the unified
felt whole of a system is set by its integration. Consciousness is integrated experience, and its degree tracks
the degree of integration, with no further ingredient required.
\end{proposition}

\subsection{A matter of degree}

If every vibe is a speck of experience, there is no sharp line at which consciousness turns on. There is only more
or less integration, so on this reading everything is conscious to some degree, a position close to panpsychism
\cite{strawson2006, goff2017}. A rock is barely integrated, a faint and diffuse near-heap of specks. A cell is a
self-maintaining unit, integrated enough to count as one thing. A brain is enormously integrated through its
fast-mixing hub. A second factor sits on top of integration, namely self-reflection, whether the integrated whole
models itself, and recursion lets it model that modeling. That is the step from being a unified experience to being
aware that one is. The ladder of consciousness is the ladder of integration plus self-reflection, the same tower as
the climb of the previous chapter. The lineage of process philosophy that takes experience as basic to every
actual occasion is the natural ancestor of this picture \cite{whitehead1929}.

\subsection{Inverting the hard problem}

A physics-first theory faces the hard problem, the question of why any arrangement of insentient matter should feel
like anything at all, with no derivation available from the physical to the felt \cite{chalmers1995}. This theory
does not solve that problem. It inverts it, so that the problem does not arise in the first place. By premise the
ground is already felt, so there is no gap to cross from matter to experience. Matter is a pattern within
experience, not a source of it. What remains is the combination problem familiar to panpsychism, the question of
how many small experiences combine into one larger experience, and that is answered by the same integration
measure, since many vibes integrate into one self whose unified experience exceeds the parts.

\begin{principle}[Inverting, not solving, the hard problem]
Because experience is taken as the ground by premise, there is no gap from insentient matter to feeling, so the
hard problem does not arise. What remains is the combination problem, addressed by the integration measure. The
theory claims the inversion and the mechanism, not a proof that the felt and the structural are identical.
\end{principle}

The real limit is worth restating. The inversion buys the disappearance of the hard problem at the cost of a
premise, and the integration measure addresses combination rather than existence. The theory does not prove that
the felt and the structural are the same thing. It assumes the felt at the base and shows how structure organizes
it.

\subsection{The variety of experience}

A natural worry is how just three tones could become the whole of experience. The answer is that the tone carries
only valence, a one-dimensional good-or-bad axis running through pain, peace, and pleasure, and all the variety of
experience lives in structure rather than in the tone. Alphabet poverty is the point. Twenty-six letters give all
of literature, and three tones give all of qualia, because the richness is in the arrangement and not in the
symbols. The structure is specified by a small set of dials.

\begin{table}[ht]
\centering
\begin{tabular}{ll}
\toprule
Dial & What it sets \\
\midrule
Modality / location & which sub-mesh, where in it \\
Pattern & the arrangement within the modality \\
Intensity & how strong \\
Valence & the good-or-bad tone, pain to peace to pleasure \\
Dynamics & how it moves and changes \\
Binding & how tightly it is bound into one \\
Level & how high on the tower \\
\bottomrule
\end{tabular}
\caption{The seven dials of qualia. Three tones fix valence, the dials fix the rest.}
\end{table}

Reading the dials gives the familiar kinds of experience. A sensation is low-level, single-modality,
pattern-rich, and lightly valenced. An emotion is high-level, multi-subsystem, and heavily valenced. A thought is
high-level, abstract, structure-rich, and lightly valenced. Each sense is a sub-mesh of a different shape, sight a
two-dimensional sheet, hearing a one-dimensional tonotopic line, smell a high-dimensional unordered space, which
predicts that smell should be vivid yet hard to name. So three tones no more limit the variety of experience than
two bits limit the variety of files. The valence is ternary, the quality is unbounded structure. This specifies
the structure of experience and not the felt-ness itself, and the combination problem stays open.

\subsection{Relation to integrated information theory}

This picture shares two ideas with integrated information theory, that integration is the quantity that matters
and that consciousness is a matter of degree \cite{tononi2004, hoel2013}. The footing differs sharply. That theory
starts from physical systems and asks how much integrated information they carry, treating the existence of
experience as the thing to explain. This theory starts from experience as the ground and uses integration to
address the combination problem rather than the existence of experience. The measure is similar in spirit, the
order of explanation is reversed.

Read all the way up, the whole substrate at its fullest integration is the largest and most unified self, so on
this reading the universe itself is conscious to the highest degree its integration allows. This is interpretation
built on the premise, not a measured result, and it is marked as such.


The universe is, most literally, a mind, a structure of consciousness in which
matter, space, and the forces are shapes of feeling. It is the boldest claim in the theory, and it is founded on a
premise that is named plainly as a premise. Experience is fundamental, integration is measured, and the
identification of consciousness with integrated experience, of panpsychism by degree, and of a conscious universe
are interpretations resting on the first of these. Said without ornament, the felt-ness is assumed and the
organization of it is what the theory describes.
